\documentclass[12pt]{article}
\usepackage[margin=0.8in]{geometry}
\usepackage{amsmath,amssymb,amsthm}
\usepackage{enumitem}
\usepackage{longtable}
\usepackage{array}
\usepackage{xcolor}
\usepackage{hyperref}
\usepackage{fancyhdr}
\usepackage{titlesec}
\usepackage{tcolorbox}
\usepackage{multicol}
\usepackage{booktabs}
\hypersetup{colorlinks=true, linkcolor=blue!50!black, urlcolor=blue!50!black}
\pagestyle{fancy}
\fancyhf{}
\rhead{\thepage}
\lhead{Engineering Math to Pure Math}
\setlength{\parskip}{0.55em}
\setlength{\parindent}{0pt}
\setlist[itemize]{topsep=0.2em,itemsep=0.2em}
\setlist[enumerate]{topsep=0.2em,itemsep=0.25em}
\newtcolorbox{focusbox}{colback=blue!3!white,colframe=blue!50!black,title=Focus}
\newtcolorbox{intuitionbox}{colback=green!3!white,colframe=green!40!black,title=Intuition}
\newtcolorbox{formalbox}{colback=gray!4!white,colframe=gray!70!black,title=Formal}
\newtcolorbox{walkbox}{colback=orange!5!white,colframe=orange!60!black,title=Walking question}
\newcommand{\N}{\mathbb{N}}
\newcommand{\Z}{\mathbb{Z}}
\newcommand{\Q}{\mathbb{Q}}
\newcommand{\R}{\mathbb{R}}
\newcommand{\C}{\mathbb{C}}
\newcommand{\eps}{\varepsilon}
\newcommand{\pow}{\mathcal{P}}
\theoremstyle{definition}
\newtheorem{definition}{Definition}
\newtheorem{example}{Example}
\newtheorem{exercise}{Exercise}
\theoremstyle{plain}
\newtheorem{theorem}{Theorem}
\newtheorem{lemma}{Lemma}
\theoremstyle{remark}
\newtheorem{remark}{Remark}

\usepackage{needspace}
\geometry{margin=0.95in}
\setlength{\parskip}{0.72em}
\setlength{\headheight}{14pt}
\newtcolorbox{trapbox}{colback=red!3!white,colframe=red!55!black,title=Common trap}
\newtcolorbox{historybox}{colback=yellow!7!white,colframe=yellow!45!black,title=History and connection}
\newtcolorbox{labbox}{colback=purple!3!white,colframe=purple!45!black,title=Proof lab}
\newcommand{\chapterline}{\vspace{0.4em}\hrule\vspace{0.8em}}
\title{\textbf{Proofs From the Ground Up}\\A Deep Conversational Reader Before Velleman}
\author{Engineering Mathematics to Pure Mathematics}
\date{May 2026}
\begin{document}
\maketitle
\tableofcontents
\newpage

\section*{How This Book Should Be Read}
\addcontentsline{toc}{section}{How This Book Should Be Read}

This is not a survey. A survey flies over the land and names the rivers. This reader is meant to put your boots in the mud. We will spend a long time on the first ideas because the first ideas carry everything later. If statements, quantifiers, definitions, sets, and functions are fuzzy, then real analysis feels like weather, abstract algebra feels like arbitrary rules, and topology feels like someone replaced geometry with riddles.

So we will not rush.

The path is the same one you asked for: from basic logic and notation toward the proof-writing maturity needed for Velleman, then onward toward analysis, algebra, topology, differential geometry, and algebraic geometry. But the emphasis here is on depth before motion. Each chapter tries to give the same concept from several angles:

\begin{itemize}
\item the informal picture,
\item the formal definition,
\item examples and non-examples,
\item what the notation is really saying,
\item the common beginner mistakes,
\item historical and later-field connections,
\item worked proofs,
\item exercises you can write,
\item questions you can think through while walking.
\end{itemize}

Mathematics is not made of symbols. It is made of ideas disciplined by symbols. The discipline matters. So does the idea.

\begin{focusbox}
Do not measure progress by pages finished. Measure it by whether you can close the PDF and explain the idea without pretending.
\end{focusbox}

\subsection*{The rule for this reader}

When a definition appears, do not merely nod. Work it. Produce examples. Break it. Repair it. Ask why every word is present. Good mathematicians do this almost automatically. Beginners often think they are supposed to skip it and get to the theorem. That is backwards. The theorem is built out of definitions.

\subsection*{What ``beginner friendly'' means here}

Beginner friendly does not mean vague. It means no step is treated as beneath explanation. We will be conversational, but we will end in exact statements. The point is not to make rigorous mathematics soft. The point is to let rigor arrive with enough handles that you can hold it.

\newpage
\section{The Change in Game: From Getting Answers to Giving Reasons}
\chapterline

Engineering mathematics often trains you to move quickly from problem type to method. That training is valuable. If you see a linear system, you think matrices. If you see a differential equation, you scan for separation, integrating factors, characteristic equations, transforms. If you see an integral, you try substitution, integration by parts, symmetry, or numerical approximation.

Pure mathematics asks you to slow down at a different point. Before asking ``what method solves this?'', it asks:

\begin{quote}
What exactly is being claimed, and why does it follow from the assumptions?
\end{quote}

That question sounds modest. It is not. It changes what counts as success.

\subsection{A calculation can be right for the wrong reason}

Suppose you notice:
\[
1=1^2,\quad 1+3=4=2^2,\quad 1+3+5=9=3^2,\quad 1+3+5+7=16=4^2.
\]

It is natural to guess:
\[
1+3+5+\cdots+(2n-1)=n^2.
\]

You can check this for $n=5,6,7,100$. You can even write a program to check it for the first million values. That is useful evidence. It is not proof.

Why not? Because the claim is not about the first million natural numbers. It is about every natural number. A proof is the device that crosses that gap.

\begin{intuitionbox}
A computation samples the landscape. A proof explains the geology.
\end{intuitionbox}

\subsection{What proof adds}

A proof gives a reason the statement cannot fail. For the odd-number identity, one proof is visual: build a square by adding L-shaped borders of size $1,3,5,\ldots$. Another proof is induction. Another uses the arithmetic series formula. The same theorem can have several proofs because the theorem has several inner explanations.

This matters for mathematical maturity. Mature readers do not just ask whether a proof is correct. They ask what the proof reveals.

\subsection{From formulas to structures}

Engineering math often focuses on formulas:
\[
F=ma,\qquad e^{i\theta}=\cos\theta+i\sin\theta,\qquad A\vec{x}=\vec{b}.
\]

Pure math often asks about structures:

\begin{itemize}
\item What is a function?
\item What is a set?
\item What is an operation?
\item What is a space?
\item What does it mean for two structures to be the same?
\item What is preserved by a map?
\end{itemize}

This is not abstraction for sport. It is compression. One theorem about groups can explain clock arithmetic, symmetries of a square, permutations, and invertible matrices. One theorem about compactness can explain why continuous functions behave well on closed intervals and why certain optimization problems have solutions.

\begin{historybox}
Many abstractions were invented after concrete methods succeeded too well. Calculus worked before its foundations were clean. Algebra solved equations before it became the study of structure. Geometry measured shapes before topology asked what ``same shape'' could mean without rulers. The rigor came because the ideas were powerful enough to need a stronger language.
\end{historybox}

\subsection{Exercise: three kinds of understanding}

For the identity
\[
1+3+\cdots+(2n-1)=n^2,
\]
write three explanations:

\begin{enumerate}
\item a computational check for $n=5$,
\item an informal visual explanation,
\item a proof plan for induction.
\end{enumerate}

Do not worry yet about perfect formal style. Notice how each explanation gives a different kind of confidence.

\begin{walkbox}
Think of a formula you know well from engineering. Ask: do I know why it is true, where it fails, and which assumptions make it work?
\end{walkbox}

\newpage
\section{Mathematical Sentences: Statements, Predicates, and Claims}
\chapterline

Before proof, there is grammar. Not school grammar, but mathematical grammar: what kind of sentence is this, what objects does it mention, and can it be true or false?

\begin{definition}
A \emph{statement} is a sentence that is either true or false.
\end{definition}

Examples:

\begin{itemize}
\item ``$7$ is prime'' is a true statement.
\item ``$9$ is prime'' is a false statement.
\item ``There are infinitely many primes'' is a true statement.
\item ``Every continuous function is differentiable'' is a false statement.
\end{itemize}

Notice that a statement does not have to be easy. For centuries, Fermat's Last Theorem was a statement whose truth was unknown. It was still a statement.

\subsection{Not every mathematical-looking phrase is a statement}

The expression
\[
x^2+1
\]
is not a statement. It does not assert anything.

The sentence
\[
x^2+1=10
\]
is not a statement until $x$ has a value or a quantifier. It is a predicate.

The command

\begin{quote}
Solve $x^2+1=10$.
\end{quote}

is not a statement either. It is an instruction.

The question

\begin{quote}
Is $x^2+1=10$?
\end{quote}

is not a statement. It asks for one.

\begin{definition}
A \emph{predicate} is a sentence with one or more variables that becomes a statement when the variables are assigned values from a specified domain.
\end{definition}

\subsection{Predicates are like Boolean-valued functions, with a warning}

In software, you might write:

\begin{verbatim}
function isPrime(n) {
  ...
}
\end{verbatim}

The expression \texttt{isPrime(n)} is not a value until \texttt{n} is supplied. In mathematics, $P(n)$ behaves similarly. But mathematics forces you to say what kind of input $n$ is allowed to be. Natural number? Integer? Real? Polynomial? Matrix?

The predicate
\[
P(x): x^2<2
\]
has different flavor depending on the domain:

\begin{itemize}
\item over $\N$, it is true only for $x=1$ if $\N=\{1,2,\ldots\}$,
\item over $\Z$, it is true for $x=-1,0,1$,
\item over $\R$, it is true for $-\sqrt2<x<\sqrt2$.
\end{itemize}

The same symbols do not always describe the same mathematical situation.

\subsection{The beginner mistake: reading predicates as statements}

A common early error is to see something like
\[
x^2\geq0
\]
and treat it as automatically true. But true where? If $x\in\R$, yes. If $x\in\C$, the symbol $\geq$ is not even naturally defined for complex numbers in the same way. The domain matters.

\begin{trapbox}
Never let a variable float without a domain. A variable without a domain is a character without a story.
\end{trapbox}

\subsection{Claims}

A theorem, proposition, lemma, or corollary is a statement that is intended to be proved. The names mostly indicate role and importance.

\begin{itemize}
\item A \emph{theorem} is a major result.
\item A \emph{proposition} is a result of moderate importance.
\item A \emph{lemma} is a helper result used to prove something else.
\item A \emph{corollary} follows quickly from a previous result.
\end{itemize}

These are cultural labels, not rigid species. One book's lemma may be another book's theorem.

\subsection{Exercises}

Classify each as statement, predicate, expression, command, or question.

\begin{enumerate}
\item $3+4=7$.
\item $x+4=7$.
\item Find all $x$ such that $x+4=7$.
\item $\sin^2 x+\cos^2 x$.
\item Every finite group of prime order is cyclic.
\item Is every ring a field?
\item $A\subseteq B$.
\item Let $f:\R\to\R$ be continuous.
\end{enumerate}

For every predicate, state at least two possible domains and explain how the meaning changes.

\newpage
\section{Variables Have Roles: Fixed, Arbitrary, Chosen, Bound}
\chapterline

Beginners often treat variables as generic letters. In proofs, variables have jobs. Confusing those jobs is one of the fastest ways to get lost.

\subsection{Fixed but unknown}

When a proof says:

\begin{quote}
Let $n\in\Z$.
\end{quote}

we introduce one integer $n$. We may not know whether it is $5$ or $-20$ or $0$, but for the duration of the argument it is fixed.

This is like receiving an input to a function. You do not know the value, but the value is fixed during that run.

\subsection{Arbitrary}

When a proof says:

\begin{quote}
Let $n\in\Z$ be arbitrary.
\end{quote}

the word arbitrary means we are not allowed to use any special property of $n$ beyond $n\in\Z$. If the proof works for this arbitrary $n$, it works for all integers.

This is how universal statements are proved.

\subsection{Chosen as a witness}

When proving
\[
\exists n\in\Z\text{ such that }n^2=49,
\]
you may choose $n=7$. That chosen object is called a witness.

Existential proofs often turn on finding the right witness. Sometimes the witness is explicit, like $7$. Sometimes it is constructed from earlier information. Sometimes existence is proved indirectly, which can be unsettling at first.

\subsection{Bound variables}

In the statement
\[
\forall x\in\R,\ x^2\geq0,
\]
the variable $x$ is bound by the quantifier $\forall$. It lives inside that sentence.

In
\[
\sum_{i=1}^n i^2,
\]
the variable $i$ is bound by the summation. It is a local variable.

This is scope. If you have programmed, you already know the smell of a scope error. Mathematics has them too.

\subsection{Free variables}

In
\[
x^2\geq0,
\]
the variable $x$ is free unless earlier context has defined it.

Free variables are not bad, but they require context. A theorem statement should not leave important variables floating ambiguously.

\subsection{Worked parsing}

Consider:
\[
\forall x\in\R,\ \exists y\in\R\text{ such that } y>x.
\]

The variable $x$ is chosen arbitrarily from $\R$. After $x$ is known, the statement promises that some $y$ can be found. A proof can say:

\begin{quote}
Let $x\in\R$ be arbitrary. Choose $y=x+1$. Then $y\in\R$ and $y>x$.
\end{quote}

Here $x$ is arbitrary, while $y$ is a witness depending on $x$.

\begin{trapbox}
Do not choose a witness before the quantifiers allow it. In $\forall x\,\exists y$, $y$ may depend on $x$. In $\exists y\,\forall x$, it may not.
\end{trapbox}

\subsection{Exercises}

For each statement, identify which variables are arbitrary and which are witnesses.

\begin{enumerate}
\item $\forall n\in\Z,\exists m\in\Z$ such that $n=2m$.
\item $\exists m\in\Z$ such that $\forall n\in\Z,\ n=2m$.
\item $\forall \eps>0,\exists N\in\N,\forall n\geq N,\ |a_n-L|<\eps$.
\item $\forall a\in A,\forall b\in B,\exists c\in C$ such that $F(a,b,c)$.
\end{enumerate}

\begin{walkbox}
When you hear a theorem, ask: which objects does the world hand me, and which objects am I allowed to choose in response?
\end{walkbox}

\newpage
\section{Quantifiers: The Game of For Every and There Exists}
\chapterline

The symbols $\forall$ and $\exists$ are small, but they decide the shape of the proof.

\[
\forall = \text{for every},\qquad \exists = \text{there exists}.
\]

\subsection{Universal statements}

A universal statement has the form:
\[
\forall x\in D,\ P(x).
\]

It says every object in the domain $D$ satisfies property $P$.

To prove it:

\begin{quote}
Let $x\in D$ be arbitrary. Show $P(x)$.
\end{quote}

To disprove it:

\begin{quote}
Find one $x\in D$ such that $P(x)$ is false.
\end{quote}

\subsection{Existential statements}

An existential statement has the form:
\[
\exists x\in D\text{ such that }P(x).
\]

It says at least one object in $D$ satisfies $P$.

To prove it:

\begin{quote}
Produce a witness $x\in D$ and verify $P(x)$.
\end{quote}

To disprove it:

\begin{quote}
Show every $x\in D$ fails $P(x)$.
\end{quote}

\subsection{The quantifier game}

Think of $\forall$ as the challenger choosing an object, and $\exists$ as you responding with an object.

Statement:
\[
\forall x\in\R,\exists y\in\R\text{ such that }x+y=0.
\]

The challenger chooses $x$. You respond with $y=-x$. You win.

Statement:
\[
\exists y\in\R,\forall x\in\R\text{ such that }x+y=0.
\]

You must choose $y$ first. Then the challenger chooses $x$. No single $y$ can cancel every $x$. You lose.

\subsection{A physical analogy}

Imagine a thermostat specification:

\begin{quote}
For every room temperature target, there exists a control setting that reaches it.
\end{quote}

That is plausible. The setting can depend on the target.

Now compare:

\begin{quote}
There exists one control setting that reaches every room temperature target.
\end{quote}

That is absurd. One setting cannot make the room both $18^\circ$C and $24^\circ$C at the same time.

\subsection{Nested quantifier practice}

Statement:
\[
\forall x\in\R,\exists y\in\R,\ y>x.
\]

True. Given $x$, choose $y=x+1$.

Statement:
\[
\exists y\in\R,\forall x\in\R,\ y>x.
\]

False. It asks for a real number larger than every real number.

Statement:
\[
\forall x\in\R,\exists y\in\R,\ y^2=x.
\]

False over $\R$, because negative $x$ has no real square root.

True if the domain of $x$ is $[0,\infty)$.

\begin{focusbox}
The domain and the quantifier order often matter more than the formula.
\end{focusbox}

\subsection{Exercises}

Decide true or false. Prove or give a counterexample.

\begin{enumerate}
\item $\forall x\in\R,\exists y\in\R,\ x<y$.
\item $\exists y\in\R,\forall x\in\R,\ x<y$.
\item $\forall n\in\Z,\exists m\in\Z,\ n<m$.
\item $\exists m\in\Z,\forall n\in\Z,\ n<m$.
\item $\forall x\in\R,\exists y\in\R,\ x+y=10$.
\item $\exists y\in\R,\forall x\in\R,\ x+y=10$.
\end{enumerate}

\newpage
\section{Negation: How to Disagree Precisely}
\chapterline

Negation is not a mood. It is a precise operation. To negate a mathematical statement, you must say exactly what it means for the statement to fail.

\subsection{Negating universal statements}

The negation of
\[
\forall x\in D,\ P(x)
\]
is
\[
\exists x\in D\text{ such that }\neg P(x).
\]

If not every object has property $P$, then at least one object fails property $P$.

\subsection{Negating existential statements}

The negation of
\[
\exists x\in D\text{ such that }P(x)
\]
is
\[
\forall x\in D,\ \neg P(x).
\]

If no object with property $P$ exists, then every object fails $P$.

\subsection{Everyday versions}

Original:

\begin{quote}
Every lock in the building opens with this key.
\end{quote}

Negation:

\begin{quote}
There is at least one lock in the building that does not open with this key.
\end{quote}

Original:

\begin{quote}
Some server is reachable.
\end{quote}

Negation:

\begin{quote}
No server is reachable.
\end{quote}

\subsection{Negating implication}

The statement
\[
P\Rightarrow Q
\]
fails exactly when $P$ is true and $Q$ is false.

So:
\[
\neg(P\Rightarrow Q)\Longleftrightarrow P\land\neg Q.
\]

This is important for counterexamples. To disprove:

\begin{quote}
If a function is continuous, then it is differentiable.
\end{quote}

you need a function that is continuous and not differentiable. The classic example is $f(x)=|x|$ at $0$.

\subsection{Negating a limit}

Definition:
\[
a_n\to L
\]
means
\[
\forall \eps>0,\exists N\in\N,\forall n\geq N,\ |a_n-L|<\eps.
\]

Negation:
\[
\exists \eps>0,\forall N\in\N,\exists n\geq N,\ |a_n-L|\geq\eps.
\]

In words: there is some tolerance band around $L$ that the sequence keeps escaping, no matter how far out you go.

\begin{historybox}
This kind of quantified precision was part of the nineteenth-century repair of calculus. Mathematicians like Cauchy and Weierstrass wanted to replace hand-wavy infinitesimals and geometric intuition with definitions that could survive pathological examples.
\end{historybox}

\subsection{Exercises}

Negate each statement in symbols and plain English.

\begin{enumerate}
\item $\forall n\in\Z,\ n^2\geq0$.
\item $\exists x\in\R,\ x^2+1=0$.
\item $\forall x\in A,\exists y\in B,\ R(x,y)$.
\item $\exists y\in B,\forall x\in A,\ R(x,y)$.
\item If $n$ is prime, then $n$ is odd.
\item Every bounded sequence converges.
\end{enumerate}

\begin{walkbox}
When you hear a universal claim in everyday life, practice finding the exact counterexample that would refute it.
\end{walkbox}

\newpage
\section{Definitions: The Machines Inside the Words}
\chapterline

In ordinary language, definitions can be loose. In mathematics, a definition is an operating manual.

If you know the definition applies, it tells you what you may use. If you want to prove the definition applies, it tells you what you must show.

\subsection{Even and odd}

\begin{definition}
An integer $n$ is even if there exists $k\in\Z$ such that $n=2k$.
\end{definition}

\begin{definition}
An integer $n$ is odd if there exists $k\in\Z$ such that $n=2k+1$.
\end{definition}

If a proof says $n$ is even, you may write $n=2k$ for some integer $k$.

If a proof asks you to show $n$ is even, you must show $n=2k$ for some integer $k$.

\subsection{Divisibility}

\begin{definition}
For integers $a$ and $b$, we say $a$ divides $b$, written $a\mid b$, if there exists $k\in\Z$ such that $b=ak$.
\end{definition}

Notice the direction: $a\mid b$ means $b$ is a multiple of $a$. Beginners often reverse it.

\subsection{Opening and closing}

Think of a definition as a box.

Opening the box:

\begin{quote}
Since $a\mid b$, there exists $k\in\Z$ such that $b=ak$.
\end{quote}

Closing the box:

\begin{quote}
Since $b=ak$ for some $k\in\Z$, we have $a\mid b$.
\end{quote}

Most early proofs are box work.

\subsection{Worked proof: even plus even}

\begin{theorem}
The sum of two even integers is even.
\end{theorem}

\begin{proof}
Let $a,b\in\Z$ be even. Since $a$ is even, there exists $m\in\Z$ such that $a=2m$. Since $b$ is even, there exists $n\in\Z$ such that $b=2n$. Then
\[
a+b=2m+2n=2(m+n).
\]
Since $m+n\in\Z$, the number $a+b$ has the form $2k$ for an integer $k$, namely $k=m+n$. Therefore $a+b$ is even.
\end{proof}

\subsection{Worked proof: divisibility is transitive}

\begin{theorem}
If $a,b,c\in\Z$, $a\mid b$, and $b\mid c$, then $a\mid c$.
\end{theorem}

\begin{proof}
Suppose $a\mid b$ and $b\mid c$. Since $a\mid b$, there exists $m\in\Z$ such that $b=am$. Since $b\mid c$, there exists $n\in\Z$ such that $c=bn$. Substituting,
\[
c=bn=(am)n=a(mn).
\]
Since $mn\in\Z$, this shows $c$ is an integer multiple of $a$. Therefore $a\mid c$.
\end{proof}

\begin{trapbox}
Do not write ``$a\mid b$, so $a=bk$.'' The definition says $b=ak$. The divided object is the one being expressed as a multiple.
\end{trapbox}

\subsection{Exercises}

Prove each directly from definitions.

\begin{enumerate}
\item The product of two even integers is even.
\item The product of two odd integers is odd.
\item If $a\mid b$ and $a\mid c$, then $a\mid(b+c)$.
\item If $a\mid b$, then $a\mid bc$ for every integer $c$.
\item If $a$ and $b$ are rational, then $a+b$ is rational.
\end{enumerate}

\newpage
\section{If-Then Statements and Proof Strategy}
\chapterline

Many theorems are implications:
\[
P\Rightarrow Q.
\]

The part $P$ is the hypothesis. The part $Q$ is the conclusion.

The proof task is not to show $P$ is true. It is to show that if $P$ is true, then $Q$ must be true.

\subsection{Direct proof}

Use direct proof when the hypothesis naturally unfolds into the conclusion.

\begin{theorem}
If $n$ is divisible by $12$, then $n$ is divisible by $3$.
\end{theorem}

\begin{proof}
Suppose $n$ is divisible by $12$. Then $n=12k$ for some $k\in\Z$. Since $12k=3(4k)$ and $4k\in\Z$, we have $n=3m$ for the integer $m=4k$. Therefore $n$ is divisible by $3$.
\end{proof}

\subsection{Contrapositive}

The contrapositive of
\[
P\Rightarrow Q
\]
is
\[
\neg Q\Rightarrow\neg P.
\]

These statements are logically equivalent.

Use the contrapositive when the negation of the conclusion is easier to work with than the hypothesis.

\begin{theorem}
If $n^2$ is even, then $n$ is even.
\end{theorem}

\begin{proof}
We prove the contrapositive: if $n$ is not even, then $n^2$ is not even. For integers, not even means odd. Suppose $n$ is odd. Then $n=2k+1$ for some $k\in\Z$. Thus
\[
n^2=(2k+1)^2=4k^2+4k+1=2(2k^2+2k)+1.
\]
Since $2k^2+2k\in\Z$, $n^2$ is odd. Hence $n^2$ is not even. Therefore, by contrapositive, if $n^2$ is even, then $n$ is even.
\end{proof}

\subsection{Contradiction}

To prove $P$ by contradiction, assume $\neg P$ and derive an impossibility.

\begin{theorem}
There is no smallest positive real number.
\end{theorem}

\begin{proof}
Suppose, for contradiction, that there is a smallest positive real number $r$. Since $r>0$, the number $r/2$ is also positive. But $r/2<r$, contradicting the claim that $r$ is the smallest positive real number. Therefore no smallest positive real number exists.
\end{proof}

\subsection{Choosing the method}

Direct proof asks: can I unfold the hypothesis and reach the conclusion?

Contrapositive asks: is the negation of the conclusion more concrete?

Contradiction asks: would denying the result trap me in an impossibility?

\begin{labbox}
For each theorem below, write only the strategy and first sentence.
\begin{enumerate}
\item If $n^2$ is odd, then $n$ is odd.
\item If $a+b$ is irrational and $a$ is rational, then $b$ is irrational.
\item There are infinitely many primes.
\item If $A\subseteq B$ and $B\subseteq C$, then $A\subseteq C$.
\end{enumerate}
\end{labbox}

\newpage
\section{Proof by Contradiction: Power and Danger}
\chapterline

Contradiction is dramatic: assume the opposite and show the world breaks. It is also easy to misuse. The goal is not to create confusion. The goal is to force a precise impossibility.

\subsection{The irrationality of $\sqrt2$}

This is one of the classic first proofs because it reveals a hidden structure in the integers: parity and divisibility constrain square roots.

\begin{theorem}
$\sqrt2$ is irrational.
\end{theorem}

\begin{proof}
Suppose, for contradiction, that $\sqrt2$ is rational. Then $\sqrt2=a/b$ for integers $a,b$ with $b\neq0$. Choose this fraction in lowest terms, so $a$ and $b$ have no common factor greater than $1$.

Squaring gives
\[
2=\frac{a^2}{b^2},
\]
so
\[
a^2=2b^2.
\]
Thus $a^2$ is even. By the theorem proved earlier, if $a^2$ is even, then $a$ is even. So $a=2k$ for some $k\in\Z$.

Substitute:
\[
(2k)^2=2b^2.
\]
Then
\[
4k^2=2b^2,
\]
so
\[
b^2=2k^2.
\]
Thus $b^2$ is even, so $b$ is even.

We have shown both $a$ and $b$ are even, so they share a factor of $2$. This contradicts the assumption that $a/b$ was in lowest terms. Therefore $\sqrt2$ is irrational.
\end{proof}

\subsection{The historical shock}

The discovery of irrational magnitudes is often associated with the Pythagorean tradition. Whether the old stories are embellished or not, the mathematical point is real: geometry produced lengths that arithmetic, if limited to ratios of whole numbers, could not contain.

The diagonal of a unit square has length $\sqrt2$. That length is perfectly visible geometrically and impossible to express as a ratio of integers. Geometry forced number to grow.

\subsection{The structure of the proof}

The proof has several moving parts:

\begin{enumerate}
\item Assume rationality.
\item Choose a fraction in lowest terms.
\item Derive that numerator and denominator are both even.
\item Contradict lowest terms.
\end{enumerate}

The contradiction is not ``something feels wrong.'' It is a direct collision between ``lowest terms'' and ``both divisible by $2$.''

\begin{trapbox}
In a contradiction proof, name the impossible thing. A vague contradiction is not enough. Say exactly which statement conflicts with which earlier assumption.
\end{trapbox}

\subsection{Exercises}

\begin{enumerate}
\item Prove that $\sqrt3$ is irrational. You may use: if $3\mid n^2$, then $3\mid n$.
\item Prove that $\sqrt5$ is irrational using the analogous divisibility fact.
\item Prove there is no rational number $r$ such that $r^2=2$.
\item Explain why the lowest-terms assumption matters.
\end{enumerate}

\newpage
\section{Examples, Non-Examples, and Counterexamples}
\chapterline

Examples are not afterthoughts. They are the laboratory where definitions become real.

\subsection{Example versus proof}

An example can prove an existential statement:

\[
\exists n\in\Z\text{ such that }n^2=25.
\]

Take $n=5$.

But an example cannot prove a universal statement:

\[
\forall n\in\Z,\ n^2\geq0.
\]

Checking $n=5$ helps you understand, but it does not prove the statement.

\subsection{Counterexamples}

A counterexample disproves a universal claim.

Claim:
\[
\forall x,y\in\R,\ x^2=y^2\Rightarrow x=y.
\]

Counterexample: $x=1$, $y=-1$. Then $x^2=y^2$, but $x\neq y$.

\subsection{Near misses}

A near miss almost satisfies a definition but fails one condition. Near misses are powerful because they reveal why the missing condition matters.

Example: A function can be injective but not surjective.

Let $f:\N\to\N$ by $f(n)=n+1$. It is injective: different inputs give different outputs. It is not surjective if $\N=\{1,2,3,\ldots\}$, because $1$ is never hit.

Example: A function can be surjective but not injective.

Let $g:\{1,2\}\to\{a\}$ by $g(1)=a$ and $g(2)=a$. It hits every element of the codomain, but merges two inputs.

\subsection{The counterexample habit}

Whenever you learn a theorem, remove one assumption and try to break it.

Theorem:

\begin{quote}
A continuous function on a closed bounded interval attains a maximum.
\end{quote}

Remove closed:
\[
f(x)=x \text{ on } (0,1)
\]
has no maximum.

Remove bounded:
\[
f(x)=x \text{ on } [0,\infty)
\]
has no maximum.

Remove continuous:
a function on $[0,1]$ that jumps upward near a missing value can fail to attain a supremum.

\begin{historybox}
Much of real analysis was shaped by counterexamples. Mathematicians learned that intuitive claims about functions could fail in strange ways. The response was not to abandon intuition, but to sharpen definitions.
\end{historybox}

\subsection{Exercises}

For each false statement, find a counterexample. Then explain which assumption fails.

\begin{enumerate}
\item Every injective function is surjective.
\item Every surjective function is injective.
\item Every bounded sequence converges.
\item Every continuous function is differentiable.
\item If $A\cup B=A\cup C$, then $B=C$.
\item If $ab=0$, then $a=0$ and $b=0$.
\end{enumerate}

\newpage
\section{Sets: The Grammar of Modern Mathematics}
\chapterline

A set is a collection of objects. The objects are called elements.

\[
x\in A
\]
means $x$ is an element of $A$.

\[
x\notin A
\]
means $x$ is not an element of $A$.

This seems almost too simple. But sets became the basic grammar of modern mathematics because almost anything can be collected: numbers, points, functions, matrices, solutions, open sets, symmetries, curves, ideals.

\subsection{Basic set operations}

\begin{definition}
The union of $A$ and $B$ is
\[
A\cup B=\{x:x\in A\text{ or }x\in B\}.
\]
\end{definition}

The ``or'' is inclusive. If $x$ belongs to both, it belongs to the union.

\begin{definition}
The intersection of $A$ and $B$ is
\[
A\cap B=\{x:x\in A\text{ and }x\in B\}.
\]
\end{definition}

\begin{definition}
The difference of $A$ and $B$ is
\[
A\setminus B=\{x:x\in A\text{ and }x\notin B\}.
\]
\end{definition}

\begin{definition}
The power set of $A$, written $\pow(A)$, is the set of all subsets of $A$.
\end{definition}

If $A=\{1,2\}$, then
\[
\pow(A)=\{\emptyset,\{1\},\{2\},\{1,2\}\}.
\]

\subsection{Subsets}

\begin{definition}
$A\subseteq B$ means every element of $A$ is an element of $B$.
\end{definition}

To prove $A\subseteq B$, use the universal structure:

\begin{quote}
Let $x\in A$ be arbitrary. Show $x\in B$.
\end{quote}

To disprove $A\subseteq B$, find one $x\in A$ with $x\notin B$.

\subsection{Set equality}

To prove $A=B$, prove:
\[
A\subseteq B\quad\text{and}\quad B\subseteq A.
\]

This is often called proving both inclusions.

\begin{intuitionbox}
Set equality is mutual containment. Nothing on the left is missing from the right, and nothing on the right is extra.
\end{intuitionbox}

\subsection{Historical connection}

Cantor's set theory made infinity into an object of study. Sets allowed mathematicians to compare the sizes of infinite collections, define functions carefully, and eventually build the foundations for analysis and topology. Set theory also produced paradoxes when handled too freely, which pushed mathematics toward more careful axioms.

\newpage
\section{Set Proofs: Element Chasing in Detail}
\chapterline

Set proofs often feel mechanical at first. That is not bad. You are learning grammar. Later, the grammar becomes invisible enough that you can think with it.

\subsection{A simple inclusion}

\begin{theorem}
For all sets $A$ and $B$, $A\cap B\subseteq A$.
\end{theorem}

\begin{proof}
Let $x\in A\cap B$ be arbitrary. By definition of intersection, $x\in A$ and $x\in B$. In particular, $x\in A$. Therefore $A\cap B\subseteq A$.
\end{proof}

\subsection{A set equality}

\begin{theorem}
For all sets $A$ and $B$,
\[
A\setminus B=A\cap B^c.
\]
\end{theorem}

\begin{proof}
Let $x$ be arbitrary. Then
\[
x\in A\setminus B
\]
if and only if $x\in A$ and $x\notin B$. The condition $x\notin B$ is equivalent to $x\in B^c$. Thus $x\in A$ and $x\in B^c$, which is equivalent to $x\in A\cap B^c$. Therefore $A\setminus B=A\cap B^c$.
\end{proof}

\subsection{A distributive law}

\begin{theorem}
For all sets $A,B,C$,
\[
A\cap(B\cup C)=(A\cap B)\cup(A\cap C).
\]
\end{theorem}

\begin{proof}
We prove both inclusions.

First let $x\in A\cap(B\cup C)$. Then $x\in A$ and $x\in B\cup C$. Since $x\in B\cup C$, either $x\in B$ or $x\in C$.

If $x\in B$, then $x\in A\cap B$, so $x\in(A\cap B)\cup(A\cap C)$.

If $x\in C$, then $x\in A\cap C$, so $x\in(A\cap B)\cup(A\cap C)$.

Thus $A\cap(B\cup C)\subseteq(A\cap B)\cup(A\cap C)$.

Now let $x\in(A\cap B)\cup(A\cap C)$. Then either $x\in A\cap B$ or $x\in A\cap C$.

If $x\in A\cap B$, then $x\in A$ and $x\in B$. Since $x\in B$, we have $x\in B\cup C$. Therefore $x\in A\cap(B\cup C)$.

If $x\in A\cap C$, then $x\in A$ and $x\in C$. Since $x\in C$, we have $x\in B\cup C$. Therefore $x\in A\cap(B\cup C)$.

Thus $(A\cap B)\cup(A\cap C)\subseteq A\cap(B\cup C)$. Therefore the sets are equal.
\end{proof}

\subsection{Exercises}

Prove each by element chasing.

\begin{enumerate}
\item $A\subseteq A\cup B$.
\item $A\cap(B\setminus C)=(A\cap B)\setminus C$.
\item $A\setminus(B\cup C)=(A\setminus B)\cap(A\setminus C)$.
\item $A\setminus(B\cap C)=(A\setminus B)\cup(A\setminus C)$.
\item If $A\subseteq B$, then $\pow(A)\subseteq\pow(B)$.
\end{enumerate}

\newpage
\section{Functions: Structured Promises, Not Just Formulas}
\chapterline

A function is one of the central objects in all of mathematics.

\begin{definition}
A function $f:A\to B$ assigns to every element of $A$ exactly one element of $B$.
\end{definition}

The set $A$ is the domain. The set $B$ is the codomain.

\subsection{Exactly one}

The phrase ``exactly one'' contains two promises:

\begin{itemize}
\item every input in $A$ gets an output,
\item no input gets two different outputs.
\end{itemize}

A relation that sends one input to two outputs is not a function. A rule that fails to define an output for some input in the domain is not a function from that domain.

\subsection{Formula is not enough}

The formula $x^2$ can define different functions:

\[
f:\R\to\R,\quad f(x)=x^2.
\]

\[
g:\R\to[0,\infty),\quad g(x)=x^2.
\]

\[
h:[0,\infty)\to[0,\infty),\quad h(x)=x^2.
\]

The formula is the same, but the functions differ because domain and codomain are part of the data.

\subsection{Image}

\begin{definition}
The image of $f:A\to B$ is
\[
\operatorname{im}(f)=\{f(a):a\in A\}.
\]
\end{definition}

The image is the part of the codomain actually hit.

For $f:\R\to\R$, $f(x)=x^2$, the codomain is $\R$, but the image is $[0,\infty)$.

\subsection{Why software engineers should care}

The type signature matters. A function's declared output type may be larger than the values it actually returns. Surjectivity asks whether the declared output space is fully covered. Injectivity asks whether identity of inputs can be recovered from outputs.

\newpage
\section{Injective, Surjective, Bijective}
\chapterline

\subsection{Injective}

\begin{definition}
A function $f:A\to B$ is injective if for all $a_1,a_2\in A$,
\[
f(a_1)=f(a_2)\Rightarrow a_1=a_2.
\]
\end{definition}

Injective means no merging. If two outputs are equal, the inputs must have been equal.

Example:
\[
f:\R\to\R,\quad f(x)=2x+1
\]
is injective. If $2a+1=2b+1$, then $a=b$.

Non-example:
\[
g:\R\to\R,\quad g(x)=x^2
\]
is not injective, since $g(1)=g(-1)$ but $1\neq-1$.

\subsection{Surjective}

\begin{definition}
A function $f:A\to B$ is surjective if for every $b\in B$, there exists $a\in A$ such that $f(a)=b$.
\end{definition}

Surjective means full coverage of the codomain.

Example:
\[
f:\R\to\R,\quad f(x)=x^3
\]
is surjective. Given $y\in\R$, choose $x=\sqrt[3]{y}$.

Non-example:
\[
g:\R\to\R,\quad g(x)=x^2
\]
is not surjective, since no real $x$ satisfies $x^2=-1$.

\subsection{Bijective}

\begin{definition}
A function is bijective if it is both injective and surjective.
\end{definition}

A bijection pairs the elements of the domain with the elements of the codomain perfectly: no merging and no missing.

\subsection{Proof patterns}

To prove injective:

\begin{quote}
Suppose $f(a_1)=f(a_2)$. Show $a_1=a_2$.
\end{quote}

To prove surjective:

\begin{quote}
Let $b\in B$ be arbitrary. Find $a\in A$ such that $f(a)=b$.
\end{quote}

\begin{labbox}
Prove that $f:\R\to\R$ given by $f(x)=5x-7$ is bijective. For injective, start with $f(a)=f(b)$. For surjective, given $y$, solve $5x-7=y$ for $x$.
\end{labbox}

\newpage
\section{Images and Inverse Images}
\chapterline

Images and inverse images are where functions meet sets.

\subsection{Image of a set}

If $f:A\to B$ and $S\subseteq A$, then
\[
f(S)=\{f(s):s\in S\}.
\]

This is the set of outputs produced by inputs from $S$.

\subsection{Inverse image of a set}

If $T\subseteq B$, then
\[
f^{-1}(T)=\{a\in A:f(a)\in T\}.
\]

This is the set of inputs whose outputs land in $T$.

\begin{trapbox}
The notation $f^{-1}(T)$ does not require $f$ to have an inverse function. It is an inverse image of a set, not necessarily an inverse function.
\end{trapbox}

\subsection{Worked proof}

\begin{theorem}
For $f:A\to B$ and $S,T\subseteq B$,
\[
f^{-1}(S\cap T)=f^{-1}(S)\cap f^{-1}(T).
\]
\end{theorem}

\begin{proof}
Let $a\in A$. Then
\[
a\in f^{-1}(S\cap T)
\]
if and only if $f(a)\in S\cap T$. This holds if and only if $f(a)\in S$ and $f(a)\in T$. That is equivalent to $a\in f^{-1}(S)$ and $a\in f^{-1}(T)$, which means $a\in f^{-1}(S)\cap f^{-1}(T)$. Therefore the sets are equal.
\end{proof}

\subsection{A place where image behaves worse}

It is always true that
\[
f(S\cap T)\subseteq f(S)\cap f(T).
\]

But equality can fail. Let $f:\R\to\R$ be $f(x)=x^2$, let $S=\{-1\}$ and $T=\{1\}$. Then $S\cap T=\emptyset$, so $f(S\cap T)=\emptyset$. But
\[
f(S)=\{1\}=f(T),
\]
so
\[
f(S)\cap f(T)=\{1\}.
\]

The function merged $-1$ and $1$. This is exactly the kind of problem injectivity controls.

\subsection{Connection to topology}

Topology defines continuity using inverse images. A function between topological spaces is continuous if the inverse image of every open set is open. This is why inverse images deserve attention now. They become a central language later.

\newpage
\section{Relations and Equivalence}
\chapterline

A relation is a way of saying which objects are connected.

Examples:

\begin{itemize}
\item equality,
\item less than,
\item divisibility,
\item congruence modulo $n$,
\item ``has the same birthday as,''
\item ``is connected by a path to.''
\end{itemize}

\subsection{Equivalence relations}

\begin{definition}
A relation $\sim$ on a set $A$ is an equivalence relation if it is reflexive, symmetric, and transitive.
\end{definition}

Reflexive:
\[
a\sim a.
\]

Symmetric:
\[
a\sim b\Rightarrow b\sim a.
\]

Transitive:
\[
a\sim b\text{ and }b\sim c\Rightarrow a\sim c.
\]

\subsection{What equivalence means}

An equivalence relation says: for the purpose at hand, these objects count as the same.

Fractions are the everyday example:
\[
\frac12,\frac24,\frac36
\]
are different symbols but represent the same rational number.

Congruence modulo $5$ says:
\[
2,7,12,17,\ldots
\]
count as the same remainder class.

\subsection{Worked proof}

\begin{theorem}
Congruence modulo $n$ is an equivalence relation on $\Z$.
\end{theorem}

\begin{proof}
Define $a\sim b$ if $n\mid(a-b)$.

Reflexive: For any $a\in\Z$, $a-a=0$, and $n\mid0$, so $a\sim a$.

Symmetric: Suppose $a\sim b$. Then $n\mid(a-b)$, so $a-b=nk$ for some $k\in\Z$. Then $b-a=-nk=n(-k)$, so $n\mid(b-a)$. Therefore $b\sim a$.

Transitive: Suppose $a\sim b$ and $b\sim c$. Then $a-b=nm$ and $b-c=n\ell$ for integers $m,\ell$. Adding gives
\[
a-c=(a-b)+(b-c)=nm+n\ell=n(m+\ell).
\]
Thus $n\mid(a-c)$, so $a\sim c$.
\end{proof}

\subsection{Quotients}

An equivalence relation partitions a set into equivalence classes. The set of equivalence classes is a quotient.

Quotients are everywhere later:

\begin{itemize}
\item modular arithmetic,
\item quotient groups,
\item quotient rings,
\item projective space,
\item gluing edges of polygons to make surfaces,
\item identifying points in topology.
\end{itemize}

\begin{walkbox}
Ask: what differences are being ignored? A quotient is often a mathematical way to say ``we no longer care about this distinction.''
\end{walkbox}

\newpage
\section{Orders: Another Kind of Relation}
\chapterline

Not every important relation is equivalence. Some relations compare rather than identify.

\begin{definition}
A relation $\leq$ on a set $A$ is a partial order if it is reflexive, antisymmetric, and transitive.
\end{definition}

Reflexive:
\[
a\leq a.
\]

Antisymmetric:
\[
a\leq b\text{ and }b\leq a\Rightarrow a=b.
\]

Transitive:
\[
a\leq b\text{ and }b\leq c\Rightarrow a\leq c.
\]

\subsection{Examples}

The usual order on $\R$ is a partial order.

Subset inclusion $\subseteq$ is a partial order on $\pow(A)$.

Divisibility is a partial order on positive integers: $a\preceq b$ if $a\mid b$.

\subsection{Partial versus total}

In a total order, any two elements can be compared. In a partial order, some pairs may be incomparable.

For subsets of $\{1,2,3\}$, the sets $\{1\}$ and $\{2\}$ are incomparable under $\subseteq$: neither is contained in the other.

\subsection{Why orders matter}

Orders appear in lattice theory, topology, logic, domain theory in computer science, divisibility, and algebraic geometry. Inclusion of ideals is a partial order. Open sets are ordered by inclusion. Dependencies in software form a partial order when cycles are absent.

\subsection{Exercises}

\begin{enumerate}
\item Prove that $\subseteq$ is a partial order on $\pow(A)$.
\item Prove that divisibility is a partial order on $\N$.
\item Give an example of a relation that is reflexive and transitive but not antisymmetric.
\item Give an example of a partial order that is not total.
\end{enumerate}

\newpage
\section{Induction: Infinite Claims, Finite Proofs}
\chapterline

Induction is a proof machine for statements indexed by natural numbers.

\begin{theorem}[Induction]
Suppose $P(n)$ is a statement for each integer $n\geq1$. If $P(1)$ is true and $P(k)\Rightarrow P(k+1)$ for every $k\geq1$, then $P(n)$ is true for all $n\geq1$.
\end{theorem}

\subsection{The domino picture}

If the first domino falls, and each domino knocks down the next, then every domino falls.

The danger is that the picture can make induction feel less precise than it is. The induction step must prove:
\[
P(k)\Rightarrow P(k+1)
\]
for an arbitrary $k$.

\subsection{Worked proof}

\begin{theorem}
For all $n\geq1$,
\[
1+2+\cdots+n=\frac{n(n+1)}2.
\]
\end{theorem}

\begin{proof}
Base case: For $n=1$, the left side is $1$ and the right side is $\frac{1(2)}2=1$, so the statement is true.

Induction step: Assume for some $k\geq1$ that
\[
1+2+\cdots+k=\frac{k(k+1)}2.
\]
We must prove
\[
1+2+\cdots+k+(k+1)=\frac{(k+1)(k+2)}2.
\]
Using the induction hypothesis,
\[
1+2+\cdots+k+(k+1)=\frac{k(k+1)}2+(k+1).
\]
Now
\[
\frac{k(k+1)}2+(k+1)
=\frac{k(k+1)+2(k+1)}2
=\frac{(k+1)(k+2)}2.
\]
So the statement holds for $k+1$. By induction, it holds for every $n\geq1$.
\end{proof}

\subsection{Common induction mistakes}

\begin{itemize}
\item Forgetting the base case.
\item Assuming $P(k+1)$ instead of $P(k)$.
\item Proving only one example of the induction step.
\item Not saying where the induction hypothesis is used.
\end{itemize}

\subsection{Exercises}

Prove by induction.

\begin{enumerate}
\item $1+3+5+\cdots+(2n-1)=n^2$.
\item $2^n\geq n+1$ for all $n\geq1$.
\item $3\mid(4^n-1)$ for all $n\geq1$.
\item $n! \geq 2^{n-1}$ for all $n\geq1$.
\end{enumerate}

\newpage
\section{Strong Induction and Recursive Thinking}
\chapterline

Strong induction lets the induction step use all previous cases, not just the immediately previous one.

\begin{theorem}[Strong induction]
Suppose $P(n)$ is a statement for integers $n\geq n_0$. If $P(n_0)$ is true and for every $k\geq n_0$, the truth of $P(n_0),P(n_0+1),\ldots,P(k)$ implies $P(k+1)$, then $P(n)$ is true for all $n\geq n_0$.
\end{theorem}

\subsection{Why it is not stronger in truth, but stronger in convenience}

Ordinary induction and strong induction are logically equivalent. But strong induction often matches the problem better. If an object can break into smaller pieces of many possible sizes, you want to assume all smaller cases are known.

\subsection{Prime factorization}

\begin{theorem}
Every integer $n\geq2$ is a product of primes.
\end{theorem}

\begin{proof}
We use strong induction.

Base case: $2$ is prime, so $2$ is a product of primes.

Induction step: Assume every integer $m$ with $2\leq m\leq k$ is a product of primes. We prove $k+1$ is a product of primes.

If $k+1$ is prime, then it is already a product of primes.

If $k+1$ is composite, then $k+1=ab$ where $2\leq a\leq k$ and $2\leq b\leq k$. By the induction hypothesis, both $a$ and $b$ are products of primes. Therefore $ab=k+1$ is a product of primes.

Thus every integer $n\geq2$ is a product of primes.
\end{proof}

\subsection{Connection to algorithms}

Strong induction resembles recursive correctness proofs. If an algorithm solves a problem by reducing it to smaller subproblems, you prove correctness by assuming the algorithm works on all smaller inputs and showing it works on the current one.

\subsection{Exercises}

\begin{enumerate}
\item Prove every integer $n\geq2$ is either prime or divisible by a prime.
\item Prove every amount of postage $n\geq12$ can be made using only $4$-cent and $5$-cent stamps.
\item Prove the Fibonacci identity $F_{n+2}=F_{n+1}+F_n$ is built into the recursive definition, then prove a nontrivial Fibonacci sum identity.
\end{enumerate}

\newpage
\section{Countability and the First Shock of Infinity}
\chapterline

Infinity is not one thing.

\begin{definition}
A set is countably infinite if its elements can be put in one-to-one correspondence with $\N$.
\end{definition}

Informally: you can list the elements in a sequence so that every element eventually appears.

\subsection{The even numbers are countable}

The map
\[
f:\N\to 2\N,\quad f(n)=2n
\]
is a bijection. So the positive even integers have the same cardinality as the positive integers.

This feels strange at first because the evens are a proper subset of the naturals. Infinite sets do that.

\subsection{The integers are countable}

List:
\[
0,1,-1,2,-2,3,-3,\ldots
\]

Every integer appears.

\subsection{The rationals are countable}

Rationals feel much larger because they are dense: between any two rationals there is another rational. But density is not cardinality.

Arrange fractions $p/q$ in a grid with $p\in\Z$ and $q\in\N$. Sweep through the grid by increasing $|p|+q$, skipping duplicates. Every rational appears after finitely many steps.

\subsection{The reals are uncountable}

Cantor's diagonal argument shows that no list of real numbers in $(0,1)$ can be complete.

Suppose someone gives a list:
\[
r_1,r_2,r_3,\ldots
\]
Write their decimal expansions. Construct a new number $s$ whose first digit differs from the first digit of $r_1$, whose second digit differs from the second digit of $r_2$, and so on. Then $s$ differs from every $r_i$ in at least one decimal place. So $s$ is not on the list.

Therefore the reals in $(0,1)$ are uncountable.

\subsection{Why this matters later}

Analysis lives in uncountable spaces. Topology studies infinite collections of open sets. Algebraic geometry studies solution sets that can be infinite but controlled by finitely many equations. Cantor's work made it possible to talk about infinity without hand waving.

\newpage
\section{Real Analysis Preview: Epsilon as a Contract}
\chapterline

Calculus says a sequence approaches a limit. Analysis asks what that sentence means.

\begin{definition}
A sequence $(a_n)$ converges to $L$ if for every $\eps>0$, there exists $N\in\N$ such that for all $n\geq N$,
\[
|a_n-L|<\eps.
\]
\end{definition}

\subsection{Plain English}

No matter how small a tolerance $\eps$ someone demands, you can go far enough into the sequence that every later term stays within that tolerance of $L$.

\subsection{Who chooses what?}

The challenger chooses $\eps>0$. You choose $N$. Then the challenger may choose any $n\geq N$, and you must be safe.

\subsection{Worked proof}

\begin{theorem}
\[
\lim_{n\to\infty}\frac1n=0.
\]
\end{theorem}

\begin{proof}
Let $\eps>0$ be arbitrary. Choose $N\in\N$ such that $N>1/\eps$. If $n\geq N$, then
\[
\left|\frac1n-0\right|=\frac1n\leq\frac1N<\eps.
\]
Therefore $1/n\to0$.
\end{proof}

\subsection{Continuity}

\begin{definition}
A function $f:\R\to\R$ is continuous at $c$ if for every $\eps>0$, there exists $\delta>0$ such that for all $x\in\R$, if $|x-c|<\delta$, then $|f(x)-f(c)|<\eps$.
\end{definition}

\begin{theorem}
$f(x)=3x+2$ is continuous at every $c\in\R$.
\end{theorem}

\begin{proof}
Let $c\in\R$ and $\eps>0$. Choose $\delta=\eps/3$. If $|x-c|<\delta$, then
\[
|f(x)-f(c)|=|(3x+2)-(3c+2)|=3|x-c|<3\delta=\eps.
\]
Thus $f$ is continuous at $c$.
\end{proof}

\subsection{Why analysis is proof training}

Analysis is quantifier training under pressure. You must track dependencies carefully: $N$ depends on $\eps$; $\delta$ depends on $\eps$ and often on the point $c$; variables must be chosen in the correct order.

\newpage
\section{Abstract Algebra Preview: Structure and Symmetry}
\chapterline

Abstract algebra begins when you stop caring what the objects are made of and start caring how they combine.

\begin{definition}
A group is a set $G$ with an operation $\cdot$ satisfying closure, associativity, identity, and inverses.
\end{definition}

\subsection{Examples}

\begin{itemize}
\item $\Z$ under addition.
\item Nonzero real numbers under multiplication.
\item Symmetries of a square under composition.
\item Invertible $n\times n$ matrices under multiplication.
\end{itemize}

\subsection{Identity is unique}

\begin{theorem}
The identity element of a group is unique.
\end{theorem}

\begin{proof}
Suppose $e$ and $e'$ are both identity elements. Since $e$ is an identity, $ee'=e'$. Since $e'$ is an identity, $ee'=e$. Therefore $e=e'$. Hence the identity is unique.
\end{proof}

\subsection{Inverses are unique}

\begin{theorem}
Every element of a group has a unique inverse.
\end{theorem}

\begin{proof}
Let $a\in G$. Suppose $b$ and $c$ are both inverses of $a$. Then $ab=e$ and $ca=e$. Now
\[
b=eb=(ca)b=c(ab)=ce=c.
\]
Therefore the inverse of $a$ is unique.
\end{proof}

\subsection{Why algebra feels abstract}

The proof did not use numbers. It used only the group laws. That is the power. Any result proved from group axioms applies to every group.

\begin{historybox}
Modern algebra grew from attempts to solve polynomial equations and understand symmetry. Galois connected the solvability of equations to groups of symmetries. Noether later helped shift the subject toward structural thinking, where operations and maps became central.
\end{historybox}

\newpage
\section{Topology Preview: Continuity Without Rulers}
\chapterline

Topology studies spaces by their open sets and continuous maps.

\begin{definition}
A topology on a set $X$ is a collection $\mathcal{T}$ of subsets of $X$ such that:
\begin{enumerate}
\item $\emptyset$ and $X$ are in $\mathcal{T}$,
\item arbitrary unions of sets in $\mathcal{T}$ are in $\mathcal{T}$,
\item finite intersections of sets in $\mathcal{T}$ are in $\mathcal{T}$.
\end{enumerate}
\end{definition}

The sets in $\mathcal{T}$ are called open sets.

\subsection{Why these axioms?}

Open sets are meant to capture neighborhoods, regions where points have room. Arbitrary unions of open regions should be open because if a point lies in one of the regions, it has room there. Finite intersections should be open because a point lying in finitely many open regions has room small enough to stay inside all of them.

Infinite intersections are not required to be open. For example:
\[
\bigcap_{n=1}^{\infty}(-1/n,1/n)=\{0\},
\]
which is not open in $\R$.

\subsection{Continuity}

\begin{definition}
If $X$ and $Y$ are topological spaces, a function $f:X\to Y$ is continuous if for every open set $U\subseteq Y$, the inverse image $f^{-1}(U)$ is open in $X$.
\end{definition}

This is why inverse images mattered earlier. Topology takes the epsilon-delta idea and expresses it without distances.

\subsection{Connection to geometry}

Topology ignores exact lengths and angles. It cares about connectedness, holes, compactness, and continuity. Differential geometry later puts smooth structure and calculus back on top of topological spaces. Algebraic geometry gives spaces built from polynomial equations a topology of their own.

\newpage
\section{Differential Geometry Preview: Calculus on Curved Worlds}
\chapterline

Differential geometry asks: how do we do calculus on spaces that are not flat?

A surface like a sphere is not globally the plane, but near any small patch it resembles the plane. This is the core idea behind manifolds.

\begin{definition}
Informally, an $n$-dimensional manifold is a space that locally looks like $\R^n$.
\end{definition}

\subsection{Charts}

A chart assigns coordinates to a small region of a manifold. Think of a map of Earth. No flat map captures the whole globe perfectly, but local maps work.

\subsection{Tangent spaces}

At each point of a smooth manifold, there is a tangent space: a linear approximation to the manifold near that point.

This connects to linear algebra. Derivatives become linear maps between tangent spaces.

\subsection{Physical connection}

Differential geometry is the language of curvature. General relativity uses curved spacetime. Robotics and computer graphics use rotations, configurations, and smooth spaces. Optimization on manifolds appears in machine learning and data analysis.

\subsection{Why proof foundations matter}

Manifolds require comfort with functions, open sets, local-to-global reasoning, equivalence of coordinate descriptions, and linear maps. The early proof habits are not separate from geometry; they are the scaffolding that lets geometry become precise.

\newpage
\section{Algebraic Geometry Preview: Equations as Shapes}
\chapterline

Algebraic geometry studies solution sets of polynomial equations using algebra.

\subsection{First examples}

The equation
\[
x^2+y^2=1
\]
defines a circle.

The equation
\[
y^2=x^3-x
\]
defines an elliptic curve.

A system of equations defines the set of points satisfying all equations simultaneously.

\begin{definition}
For polynomials $f_1,\ldots,f_m\in k[x_1,\ldots,x_n]$ over a field $k$, the affine variety they define is
\[
V(f_1,\ldots,f_m)=\{a\in k^n:f_i(a)=0\text{ for all }i\}.
\]
\end{definition}

\subsection{Ideals}

\begin{definition}
An ideal $I$ in a ring $R$ is a nonempty subset such that:
\begin{enumerate}
\item if $a,b\in I$, then $a+b\in I$,
\item if $r\in R$ and $a\in I$, then $ra\in I$.
\end{enumerate}
\end{definition}

Ideals encode equations. Geometry becomes algebra when you study the polynomials that vanish on a shape.

\subsection{Why the early topics matter}

Algebraic geometry uses:

\begin{itemize}
\item sets to describe solution collections,
\item functions to describe maps between spaces,
\item rings and ideals from algebra,
\item topology to describe geometric behavior,
\item quotients to identify algebraic structures,
\item proof maturity to survive abstraction.
\end{itemize}

\begin{historybox}
Algebraic geometry began with concrete curves and surfaces. Over time it absorbed commutative algebra, topology, category theory, and number theory. Grothendieck's revolution was not making the subject obscure for fun; it created a language powerful enough to see many older phenomena as parts of one architecture.
\end{historybox}

\newpage
\section{Proof Debugging: Turning Fog Into Work}
\chapterline

When a proof fails, do not panic and do not start over immediately. Debug it.

\subsection{The proof debugger}

\begin{enumerate}
\item State the claim exactly.
\item Identify all variables and domains.
\item Circle the quantifiers.
\item Separate hypotheses from conclusion.
\item Expand definitions in the conclusion.
\item Expand definitions in the hypotheses.
\item Choose a strategy.
\item Write a skeleton.
\item Fill in one line at a time.
\item Check every line: why does it follow?
\end{enumerate}

\subsection{Gap types}

\begin{longtable}{p{0.22\linewidth}p{0.34\linewidth}p{0.34\linewidth}}
\toprule
Gap & Symptom & Repair \\
\midrule
Quantifier gap & Object used before being introduced & Introduce it with domain and role \\
Definition gap & Term used vaguely & Replace term with definition \\
Existence gap & Object claimed without witness & Construct or cite existence \\
Implication gap & A line jumps too far & Add intermediate justification \\
Case gap & Not all possibilities handled & List cases explicitly \\
Converse error & Direction reversed & Check hypothesis and conclusion \\
\bottomrule
\end{longtable}

\subsection{Bad proof example}

Claim: If $n^2$ is even, then $n$ is even.

Bad proof:

\begin{quote}
Since $n^2$ is even, $n^2=2k$. Taking square roots gives $n=\sqrt{2k}$, so $n$ is even.
\end{quote}

The broken line is the last one. Being $\sqrt{2k}$ does not mean being twice an integer. Repair by using the contrapositive.

\subsection{What good feedback sounds like}

Bad feedback:

\begin{quote}
This proof is wrong.
\end{quote}

Useful feedback:

\begin{quote}
The first unsupported step is line 3. You need to show $\sqrt{2k}$ has the form $2m$ for some integer $m$, but that is false in general.
\end{quote}

\newpage
\section{A Serious Velleman Readiness Plan}
\chapterline

You are ready to read Velleman well when the following tasks feel possible, even if not effortless.

\subsection{Core readiness}

\begin{enumerate}
\item Translate English statements into quantifiers.
\item Negate nested quantified statements.
\item Prove direct claims from definitions.
\item Use contrapositive when the conclusion has a useful negation.
\item Use contradiction with a named contradiction.
\item Prove set equality by two inclusions.
\item Prove injectivity and surjectivity from definitions.
\item Work with inverse images of sets.
\item Prove a relation is an equivalence relation.
\item Write induction proofs with base case, hypothesis, and step.
\end{enumerate}

\subsection{Four-week deep launch}

\begin{longtable}{p{0.08\linewidth}p{0.22\linewidth}p{0.34\linewidth}p{0.26\linewidth}}
\toprule
Week & Theme & Daily work & Milestone \\
\midrule
1 & Sentences and quantifiers & Translate, negate, identify variable roles & Quantifier order feels meaningful \\
2 & Definitions and proof moves & Parity, divisibility, rationality, contradiction & You can write clean one-page proofs \\
3 & Sets and functions & Element chasing, image, inverse image, injective/surjective & You can prove function claims from definitions \\
4 & Relations and induction & Equivalence, orders, induction, countability preview & You can begin Velleman without fog dominating \\
\bottomrule
\end{longtable}

\subsection{Portfolio}

Write polished proofs of:

\begin{enumerate}
\item sum of evens is even,
\item product of odds is odd,
\item divisibility is transitive,
\item if $n^2$ is even then $n$ is even,
\item $\sqrt2$ is irrational,
\item $A\cap(B\cup C)=(A\cap B)\cup(A\cap C)$,
\item $f^{-1}(S\cap T)=f^{-1}(S)\cap f^{-1}(T)$,
\item composition of injective functions is injective,
\item congruence modulo $n$ is an equivalence relation,
\item $1+2+\cdots+n=n(n+1)/2$.
\end{enumerate}

\newpage
\section{Second Pass: Logic as a Language You Can Speak}
\chapterline

The first pass taught the parts. This second pass is where the parts become usable. We revisit the same concepts, but from different angles. This is not repetition for padding. It is repetition because the same idea looks different when you approach it as a reader, a writer, a debugger, and a future analyst or algebraist.

\subsection{The reader's view}

When you read a mathematical sentence, your first job is to identify its grammar. Consider:
\[
\forall n\in\Z,\ n^2+n\text{ is even}.
\]

The outer shape is universal. It says every integer has a property. A proof must begin with an arbitrary integer:

\begin{quote}
Let $n\in\Z$ be arbitrary.
\end{quote}

The property is ``$n^2+n$ is even.'' By definition, proving this means finding an integer $k$ such that
\[
n^2+n=2k.
\]

Now the proof has a target. We are not just trying to say something plausible about $n^2+n$. We are trying to put it into the shape $2k$.

\subsection{The writer's view}

A proof might go:

\begin{proof}
Let $n\in\Z$ be arbitrary. Then
\[
n^2+n=n(n+1).
\]
One of $n$ and $n+1$ is even, so $n(n+1)$ is even. Therefore $n^2+n$ is even.
\end{proof}

This proof is short, but it hides a lemma: one of two consecutive integers is even. For a beginner portfolio, you might expand it:

\begin{proof}
Let $n\in\Z$ be arbitrary. If $n$ is even, then $n=2m$ for some $m\in\Z$, so
\[
n^2+n=n(n+1)=2m(n+1),
\]
which is even. If $n$ is odd, then $n=2m+1$ for some $m\in\Z$, so
\[
n+1=2m+2=2(m+1).
\]
Thus
\[
n(n+1)=n\cdot2(m+1)=2(n(m+1)),
\]
which is even. In both cases, $n^2+n$ is even.
\end{proof}

The longer proof is not more advanced. It is more explicit. At the beginning, explicit is good.

\subsection{The debugger's view}

Bad proof:

\begin{quote}
$n^2+n$ is even because it has two terms and one is odd and one is even.
\end{quote}

Diagnosis:

\begin{itemize}
\item ``two terms'' has no direct connection to evenness;
\item $n^2$ and $n$ are not always one odd and one even;
\item the proof never uses the definition of even.
\end{itemize}

Repair:

\begin{itemize}
\item factor $n^2+n=n(n+1)$;
\item use the fact that consecutive integers have opposite parity;
\item close the definition of even.
\end{itemize}

\subsection{The later-math view}

The expression $n(n+1)$ being even is a divisibility statement. It says:
\[
2\mid n(n+1).
\]

In number theory and abstract algebra, divisibility becomes a structural idea. In modular arithmetic, the same proof can be written:
\[
n\equiv0\pmod2\quad\text{or}\quad n\equiv1\pmod2.
\]
If $n\equiv0$, then $n(n+1)\equiv0$. If $n\equiv1$, then $n+1\equiv0$, so again $n(n+1)\equiv0$. The same idea has moved into the language of congruence.

\begin{labbox}
Write three proofs that $n^2+n$ is even: one by cases, one using modular arithmetic, and one using the phrase ``consecutive integers.'' Then decide which proof is best for a beginner and which is best for a reader who already knows congruences.
\end{labbox}

\newpage
\section{Second Pass: Quantifier Order in Real Problems}
\chapterline

Quantifier order is easy to understand in toy examples and easy to mishandle in serious mathematics. The way to improve is to translate each statement into a game of choices.

\subsection{Uniform versus pointwise}

Compare:
\[
\forall x\in A,\exists \delta>0,\ P(x,\delta)
\]
with
\[
\exists \delta>0,\forall x\in A,\ P(x,\delta).
\]

The first allows $\delta$ to depend on $x$. The second requires one $\delta$ that works for all $x$.

This distinction appears everywhere.

\subsection{Continuity versus uniform continuity}

Continuity at each point says, roughly:
\[
\forall x,\forall \eps>0,\exists \delta>0,\forall y,\text{ if }|y-x|<\delta\text{ then }|f(y)-f(x)|<\eps.
\]

The $\delta$ may depend on both $\eps$ and the point $x$.

Uniform continuity says:
\[
\forall \eps>0,\exists \delta>0,\forall x,\forall y,\text{ if }|y-x|<\delta\text{ then }|f(y)-f(x)|<\eps.
\]

Now $\delta$ depends only on $\eps$, not on the location $x$.

\subsection{Physical picture}

Imagine a road with changing steepness. Ordinary continuity says: at every location, if you zoom in enough, small horizontal movement produces small vertical movement. But the amount of zoom may depend on where you stand.

Uniform continuity says: one zoom level works across the whole road.

On a closed bounded interval, continuous functions are uniformly continuous. That theorem is a jewel of real analysis because it converts local control into global control. Compactness is the hidden machine behind that conversion.

\subsection{Software picture}

Suppose a function has a stability guarantee:

\begin{quote}
For every input, there exists a tolerance setting that keeps output error below $0.01$.
\end{quote}

That is pointwise. It may require retuning for every input.

Now compare:

\begin{quote}
There exists one tolerance setting that works for every input.
\end{quote}

That is uniform. It is much stronger and much more useful operationally.

\subsection{Exercises}

For each pair, explain which statement is stronger.

\begin{enumerate}
\item $\forall x\in A,\exists y\in B,\ R(x,y)$ versus $\exists y\in B,\forall x\in A,\ R(x,y)$.
\item $\forall p\in X,\exists r>0,\ B_r(p)\subseteq U$ versus $\exists r>0,\forall p\in X,\ B_r(p)\subseteq U$.
\item For every bug there exists a test that catches it versus there exists one test that catches every bug.
\item For every city there is a road to the capital versus there is one road that goes from every city to the capital.
\end{enumerate}

\begin{walkbox}
When walking, invent three everyday statements of the form ``for every, there exists.'' Then reverse the quantifiers and ask whether the new statement is still true.
\end{walkbox}

\newpage
\section{Second Pass: Definitions Under a Microscope}
\chapterline

Definitions are compact. That is why they are useful and why they are dangerous. A definition may compress a page of intuition into one line. If you read only the line, you may miss the intuition. If you keep only the intuition, you may lose the rigor.

\subsection{Case study: injective}

Formal definition:
\[
f:A\to B\text{ is injective if }\forall a_1,a_2\in A,\ f(a_1)=f(a_2)\Rightarrow a_1=a_2.
\]

Plain version: equal outputs force equal inputs.

Visual version: no two distinct inputs point to the same output.

Information version: the output retains enough information to recover the input.

Proof version: to prove injective, assume $f(a_1)=f(a_2)$ and show $a_1=a_2$.

Counterexample version: to disprove injective, find $a_1\neq a_2$ with $f(a_1)=f(a_2)$.

\subsection{Why the official definition uses equal outputs}

Beginners often define injective as:

\begin{quote}
Different inputs give different outputs.
\end{quote}

That is true:
\[
a_1\neq a_2\Rightarrow f(a_1)\neq f(a_2).
\]

But the standard definition uses the contrapositive:
\[
f(a_1)=f(a_2)\Rightarrow a_1=a_2.
\]

Why? Because in algebraic proofs, it is often easier to start with equality of outputs and manipulate equations until equality of inputs appears.

\subsection{Case study: open set}

Metric definition:

\[
U\subseteq X\text{ is open if }\forall x\in U,\exists r>0,\ B_r(x)\subseteq U.
\]

Plain version: every point of $U$ has wiggle room inside $U$.

Physical version: if you stand at any point in the region, you can move a little in any direction without leaving the region.

Proof version: to prove open, start with arbitrary $x\in U$ and produce a radius $r>0$.

Counterexample version: to prove not open, find a point in the set with no positive-radius ball contained in the set.

\subsection{The boundary tells the story}

The interval $(0,1)$ is open in $\R$. If $x\in(0,1)$, choose
\[
r=\frac12\min(x,1-x).
\]
Then the ball $(x-r,x+r)$ stays inside $(0,1)$.

The interval $[0,1]$ is not open in $\R$ because at $x=0$, every ball around $0$ contains negative numbers outside $[0,1]$.

\subsection{Later connection}

Open sets become the basic language of topology. The same definition pattern appears in manifolds: local neighborhoods look like open subsets of $\R^n$. Algebraic geometry has its own topologies, where open sets may behave very differently from metric open sets.

\begin{labbox}
For each definition you meet this week, write six versions: formal, plain, visual, proof, counterexample, and later connection.
\end{labbox}

\newpage
\section{Second Pass: Proofs as Programs With Human Readers}
\chapterline

A proof resembles a program in that each line must be justified. But a proof is also unlike a program because its compiler is a human reader with mathematical expectations.

\subsection{Types}

Every object has a type:

\begin{itemize}
\item $n\in\Z$,
\item $x\in\R$,
\item $A\subseteq X$,
\item $f:A\to B$,
\item $G$ is a group,
\item $U$ is open.
\end{itemize}

Type errors in proofs are common. You cannot write $x\in f$ if $f$ is a function and $x$ is not a function element. You cannot ask whether $A< B$ if $A$ and $B$ are sets without an order relation.

\subsection{Preconditions and postconditions}

In a theorem:

\begin{quote}
If $f:A\to B$ is injective and $g:B\to C$ is injective, then $g\circ f:A\to C$ is injective.
\end{quote}

The preconditions are:

\begin{itemize}
\item $f$ is a function from $A$ to $B$,
\item $g$ is a function from $B$ to $C$,
\item $f$ is injective,
\item $g$ is injective.
\end{itemize}

The postcondition is:

\begin{itemize}
\item $g\circ f$ is injective.
\end{itemize}

Proof writing becomes easier when you treat the conclusion as a spec. What exactly must be shown for $g\circ f$ to satisfy the injective spec?

\subsection{Invariants}

In induction, the statement $P(n)$ is like an invariant over natural numbers. The base case initializes it. The induction step preserves it from $k$ to $k+1$.

This analogy is not perfect, but it helps. A loop invariant must hold before and after each iteration. An induction statement must propagate from one case to the next.

\subsection{Readable proof style}

A proof should not be cryptic. Good proof prose says what is happening:

\begin{quote}
We prove both inclusions.
\end{quote}

\begin{quote}
Let $x\in A$ be arbitrary.
\end{quote}

\begin{quote}
By injectivity of $g$, we get $f(a_1)=f(a_2)$.
\end{quote}

These phrases are not filler. They orient the reader.

\subsection{Exercise}

Take a proof you wrote. Annotate it like code:

\begin{itemize}
\item mark variable introductions,
\item mark definitions opened,
\item mark definitions closed,
\item mark theorem or lemma calls,
\item mark the exact line where the conclusion is reached.
\end{itemize}

\newpage
\section{Second Pass: Set Algebra and Boolean Thinking}
\chapterline

Set operations mirror logical operations.

\[
x\in A\cap B\quad\Longleftrightarrow\quad x\in A\land x\in B.
\]

\[
x\in A\cup B\quad\Longleftrightarrow\quad x\in A\lor x\in B.
\]

\[
x\in A^c\quad\Longleftrightarrow\quad \neg(x\in A).
\]

This is why set identities often resemble logic laws.

\subsection{De Morgan's laws}

\[
(A\cup B)^c=A^c\cap B^c.
\]

\[
(A\cap B)^c=A^c\cup B^c.
\]

These are set versions of:

\[
\neg(P\lor Q)\Longleftrightarrow \neg P\land\neg Q.
\]

\[
\neg(P\land Q)\Longleftrightarrow \neg P\lor\neg Q.
\]

\subsection{Worked proof}

\begin{theorem}
\[
(A\cup B)^c=A^c\cap B^c.
\]
\end{theorem}

\begin{proof}
Let $x$ be arbitrary. Then
\[
x\in(A\cup B)^c
\]
if and only if $x\notin A\cup B$. This holds if and only if it is not the case that $x\in A$ or $x\in B$. By De Morgan's law for logic, this means $x\notin A$ and $x\notin B$. Therefore $x\in A^c$ and $x\in B^c$, which means $x\in A^c\cap B^c$. Hence the sets are equal.
\end{proof}

\subsection{Why this matters}

Analysis uses set operations constantly: unions of open sets, intersections of closed sets, complements, limsup and liminf of sets. Probability uses sigma-algebras of sets. Topology is built from collections of open sets. Algebraic geometry defines closed sets as zero sets of polynomials and studies their complements as open sets.

\subsection{Exercises}

\begin{enumerate}
\item Prove $(A\cap B)^c=A^c\cup B^c$.
\item Prove $A\setminus B=A\cap B^c$.
\item Prove $A\subseteq B$ if and only if $B^c\subseteq A^c$.
\item Translate the previous item into an implication and its contrapositive.
\end{enumerate}

\newpage
\section{Second Pass: Functions as Maps of Structure}
\chapterline

Functions do more than assign values. In higher mathematics, functions preserve or reveal structure.

\subsection{Maps in different fields}

In analysis, continuous functions preserve nearness.

In algebra, homomorphisms preserve operations:
\[
\phi(ab)=\phi(a)\phi(b).
\]

In topology, homeomorphisms preserve topological structure.

In differential geometry, smooth maps preserve differentiable structure.

In algebraic geometry, morphisms preserve polynomial/algebraic structure.

The word ``map'' is not casual. Mathematics often studies objects by studying the structure-preserving maps between them.

\subsection{Kernel as what a map forgets}

In group theory, the kernel of a homomorphism $\phi:G\to H$ is
\[
\ker\phi=\{g\in G:\phi(g)=e_H\}.
\]

This set contains the elements that become invisible under $\phi$.

In linear algebra, the null space of a matrix is the set of vectors sent to zero. Same idea.

\subsection{Quotients as controlled forgetting}

If a map forgets some information, a quotient often formalizes exactly what has been forgotten. The first isomorphism theorem says, roughly:

\[
G/\ker\phi\cong \operatorname{im}\phi.
\]

Translation: once you identify elements that the map cannot distinguish, what remains is essentially the image.

This is far beyond the first proof course, but the seed is already in functions, equivalence relations, and inverse images.

\subsection{Exercise}

For each function, say what it forgets.

\begin{enumerate}
\item $f:\Z\to\{0,1\}$ where $f(n)$ is the parity of $n$.
\item $g:\R^2\to\R$ where $g(x,y)=x$.
\item $h:\R\to\R$ where $h(x)=x^2$.
\item A hash function from strings to fixed-length hashes.
\end{enumerate}

\begin{walkbox}
Think of a function you use in software. What information does it preserve? What information does it discard? If two inputs have the same output, what relation does that define?
\end{walkbox}

\newpage
\section{Second Pass: Relations, Classification, and Quotients}
\chapterline

Equivalence relations classify objects. They say which objects belong in the same bucket.

\subsection{Classification by invariant}

An invariant is a property that stays the same under some relation or transformation.

For congruence modulo $5$, the invariant is remainder after division by $5$.

For similar triangles, the invariant is shape, not size.

For topological spaces up to homeomorphism, invariants include connectedness, compactness, and number of holes.

\subsection{Equivalence classes}

If $\sim$ is an equivalence relation on $A$, the equivalence class of $a$ is
\[
[a]=\{x\in A:x\sim a\}.
\]

For congruence modulo $5$:

\[
[2]=\{\ldots,-8,-3,2,7,12,17,\ldots\}.
\]

\subsection{Partitions}

Equivalence classes partition the set:

\begin{itemize}
\item every element belongs to some class,
\item different classes are either identical or disjoint.
\end{itemize}

This is a theorem worth proving.

\begin{theorem}
If $\sim$ is an equivalence relation on $A$, then any two equivalence classes are either disjoint or equal.
\end{theorem}

\begin{proof}
Suppose $[a]$ and $[b]$ are not disjoint. Then there exists $x\in[a]\cap[b]$. Thus $x\sim a$ and $x\sim b$. By symmetry, $a\sim x$. Since $a\sim x$ and $x\sim b$, transitivity gives $a\sim b$.

We show $[a]=[b]$. Let $y\in[a]$. Then $y\sim a$. Since $a\sim b$, transitivity gives $y\sim b$, so $y\in[b]$. Thus $[a]\subseteq[b]$.

Similarly, using $b\sim a$, we get $[b]\subseteq[a]$. Therefore $[a]=[b]$.
\end{proof}

\subsection{Why classification matters}

Mathematics often classifies objects up to an equivalence relation:

\begin{itemize}
\item fractions up to equality of ratios,
\item geometric figures up to congruence or similarity,
\item groups up to isomorphism,
\item spaces up to homeomorphism,
\item manifolds up to diffeomorphism.
\end{itemize}

Classification is one of the great mathematical projects: understand all objects of a certain kind, considering two objects the same when they have the same relevant structure.

\newpage
\section{Second Pass: Induction Beyond Formulas}
\chapterline

Induction is not only for sums. It is for any structure built step by step.

\subsection{Recursive definitions}

The Fibonacci numbers are defined by:

\[
F_0=0,\quad F_1=1,\quad F_{n+2}=F_{n+1}+F_n.
\]

The definition itself is recursive. Proofs about Fibonacci numbers often use induction because the object is built inductively.

\subsection{Worked Fibonacci proof}

\begin{theorem}
For all $n\geq1$,
\[
F_1+F_2+\cdots+F_n=F_{n+2}-1.
\]
\end{theorem}

\begin{proof}
Base case: For $n=1$, the left side is $F_1=1$, and the right side is $F_3-1=2-1=1$.

Induction step: Suppose
\[
F_1+\cdots+F_k=F_{k+2}-1.
\]
Then
\[
F_1+\cdots+F_k+F_{k+1}=(F_{k+2}-1)+F_{k+1}.
\]
Using the Fibonacci recurrence,
\[
F_{k+2}+F_{k+1}=F_{k+3}.
\]
Thus
\[
F_1+\cdots+F_{k+1}=F_{k+3}-1.
\]
This is the desired statement for $k+1$. Therefore the formula holds for all $n\geq1$.
\end{proof}

\subsection{Structural induction preview}

In computer science, data structures like lists and trees are recursive. To prove something about all lists, you prove it for the empty list and prove that adding an element preserves the property. To prove something about trees, you prove it for leaves and show it is preserved when subtrees are joined.

That is structural induction. It is the same spirit as ordinary induction but applied to recursively built objects.

\subsection{Exercise}

Write an induction proof for the length of concatenated lists:

\[
\operatorname{len}(L_1\mathbin{++}L_2)=\operatorname{len}(L_1)+\operatorname{len}(L_2).
\]

State clearly what object you are inducting on.

\newpage
\section{Second Pass: Analysis as Quantifier Discipline}
\chapterline

Analysis is where vague words like ``approaches,'' ``near,'' and ``small'' are put under oath.

\subsection{The anatomy of an epsilon proof}

To prove $a_n\to L$, the structure is:

\begin{enumerate}
\item Let $\eps>0$ be arbitrary.
\item Choose $N$ depending on $\eps$.
\item Let $n\geq N$.
\item Show $|a_n-L|<\eps$.
\end{enumerate}

The proof often has a scratch-work phase where you discover how to choose $N$.

\subsection{Scratch work versus proof}

For $a_n=3/n$, we want:
\[
\left|\frac3n-0\right|<\eps.
\]

This is:
\[
\frac3n<\eps.
\]

Solving for $n$ gives:
\[
n>\frac3\eps.
\]

That scratch work suggests choosing $N>3/\eps$.

Now the proof:

\begin{proof}
Let $\eps>0$. Choose $N\in\N$ with $N>3/\eps$. If $n\geq N$, then
\[
\left|\frac3n\right|=\frac3n\leq\frac3N<\eps.
\]
Thus $3/n\to0$.
\end{proof}

\subsection{Why scratch work is not the proof}

Scratch work often runs backward from the desired conclusion. The final proof must run forward from the chosen $N$ to the conclusion.

This distinction matters in all proof writing. Discovery and presentation are not identical.

\subsection{Exercise}

Prove:

\[
\lim_{n\to\infty}\frac{5}{n+2}=0.
\]

First write scratch work. Then write the formal proof.

\newpage
\section{Second Pass: Algebra as Lawful Manipulation}
\chapterline

Abstract algebra asks what follows from laws. The objects may be numbers, matrices, symmetries, functions, or something stranger. The laws are what matter.

\subsection{Cancellation in groups}

\begin{theorem}
If $G$ is a group and $ab=ac$, then $b=c$.
\end{theorem}

\begin{proof}
Suppose $ab=ac$. Multiply on the left by $a^{-1}$:
\[
a^{-1}(ab)=a^{-1}(ac).
\]
By associativity,
\[
(a^{-1}a)b=(a^{-1}a)c.
\]
Since $a^{-1}a=e$,
\[
eb=ec.
\]
Since $e$ is the identity,
\[
b=c.
\]
\end{proof}

\subsection{Why not divide?}

In ordinary algebra, you might say ``divide by $a$.'' In a group, the operation may not be multiplication of numbers, and it may not be commutative. The correct group-theoretic action is multiplying by the inverse on the correct side.

\subsection{Noncommutativity}

In many groups, $ab\neq ba$. Matrix multiplication is the standard example. Symmetries of a square also do not always commute: rotating then reflecting can differ from reflecting then rotating.

This is why order matters in algebraic proofs.

\subsection{Exercise}

Prove right cancellation: if $ba=ca$, then $b=c$.

Then explain why your proof must multiply by $a^{-1}$ on the right rather than on the left.

\newpage
\section{Second Pass: Topology as the Study of Stable Nearness}
\chapterline

Topology begins with a bold move: stop measuring distance directly and keep only the structure needed to talk about continuity.

\subsection{Open sets as observable neighborhoods}

In $\R$, an open interval like $(2,5)$ has the property that each point inside has room to move. The point $2.1$ has a small interval around it inside $(2,5)$. The point $2$ is not inside, so it does not matter.

In $[2,5]$, the point $2$ is inside but has no room to move left while staying inside. That is why $[2,5]$ is not open in $\R$.

\subsection{Relative openness}

But $[2,5]$ can be open as a subset of itself with the subspace topology. This is a major beginner surprise: openness depends on the ambient space.

In the space $X=[2,5]$, the set $[2,3)$ is open in $X$ because it equals:
\[
X\cap(1,3).
\]

The surrounding universe changed.

\subsection{Why this matters}

Manifolds are locally modeled on open subsets of $\R^n$. Algebraic geometry uses topologies where closed sets are solution sets of polynomial equations. Functional analysis studies topologies on spaces of functions. The word ``open'' is portable because it is defined structurally.

\subsection{Exercise}

In the subspace $X=[0,1]\subseteq\R$, decide whether the following are open in $X$:

\begin{enumerate}
\item $[0,1/2)$,
\item $(0,1/2)$,
\item $(1/2,1]$,
\item $\{0\}$.
\end{enumerate}

For each open set, express it as $X\cap U$ for some open $U\subseteq\R$.

\newpage
\section{Second Pass: Geometry, Coordinates, and What Is Real}
\chapterline

Coordinates are useful, but they are not the object itself.

\subsection{The coordinate trap}

A point on Earth can be described by latitude and longitude. But the point is not the coordinate pair. Different coordinate systems can describe the same point. Near the poles, familiar coordinates behave badly.

Differential geometry takes this seriously. A manifold is studied through charts, but the manifold is not any one chart.

\subsection{Local versus global}

The Earth is locally flat enough for a city map. Globally it is not flat. This local/global tension appears everywhere:

\begin{itemize}
\item locally Euclidean spaces can have global holes,
\item local derivatives can assemble into global curvature,
\item local coordinate formulas must agree on overlaps.
\end{itemize}

\subsection{Transition maps}

If two charts describe overlapping regions of a manifold, the transition map tells how to convert coordinates from one chart to the other. Smooth manifolds require these transition maps to be smooth.

This is where functions, inverse functions, open sets, and differentiability all meet.

\subsection{Exercise}

Explain in your own words:

\begin{enumerate}
\item why one flat map cannot represent Earth perfectly,
\item why local maps are still useful,
\item how this relates to the definition of a manifold,
\item why proof-based definitions are needed to prevent coordinate confusion.
\end{enumerate}

\newpage
\section{Second Pass: Algebraic Geometry as Translation}
\chapterline

Algebraic geometry is a translation project between equations and shapes.

\subsection{From equation to shape}

The polynomial
\[
f(x,y)=x^2+y^2-1
\]
defines the circle:
\[
V(f)=\{(x,y)\in\R^2:x^2+y^2-1=0\}.
\]

\subsection{From shape to functions}

Instead of only asking which points lie on the circle, ask which polynomial functions on the plane become the same when restricted to the circle.

On the circle,
\[
x^2+y^2=1.
\]
So the polynomial $x^2+y^2$ behaves like the constant function $1$ on that shape.

This suggests quotienting the polynomial ring by the relation generated by $x^2+y^2-1$:
\[
\R[x,y]/(x^2+y^2-1).
\]

That quotient is a first glimpse of the coordinate ring.

\subsection{Why ideals are natural}

If a polynomial vanishes on a shape, then any multiple of it also vanishes on the shape. If two polynomials vanish on the shape, their sum vanishes too. That is exactly the behavior of an ideal.

\subsection{The same early tools return}

\begin{itemize}
\item Sets: solution sets.
\item Functions: polynomial maps.
\item Equivalence relations: quotient rings.
\item Ideals: algebraic encoding of equations.
\item Topology: closed sets as varieties.
\item Proofs: every translation must be justified.
\end{itemize}

\subsection{Exercise}

Let $V=V(y-x^2)\subseteq\R^2$, the parabola.

\begin{enumerate}
\item Explain why $y$ and $x^2$ define the same function on $V$.
\item Explain why $y-x^2$ vanishes on $V$.
\item Give another polynomial that vanishes on $V$.
\item What information does the quotient $\R[x,y]/(y-x^2)$ seem to remember?
\end{enumerate}

\newpage
\section{Guided Solutions and Commentary}
\chapterline

This section is not for passive reading. Try the problems first. Then use these as models for how solutions breathe.

\subsection{Solution: sum of two odd integers}

Claim: The sum of two odd integers is even.

\begin{proof}
Let $a$ and $b$ be odd integers. Then there exist $m,n\in\Z$ such that $a=2m+1$ and $b=2n+1$. Therefore
\[
a+b=(2m+1)+(2n+1)=2m+2n+2=2(m+n+1).
\]
Since $m+n+1\in\Z$, the sum $a+b$ is even.
\end{proof}

Commentary: The proof opens the definition of odd for both numbers and closes the definition of even at the end.

\subsection{Solution: composition of surjections}

Claim: If $f:A\to B$ and $g:B\to C$ are surjective, then $g\circ f:A\to C$ is surjective.

\begin{proof}
Let $c\in C$ be arbitrary. Since $g$ is surjective, there exists $b\in B$ such that $g(b)=c$. Since $f$ is surjective, there exists $a\in A$ such that $f(a)=b$. Then
\[
(g\circ f)(a)=g(f(a))=g(b)=c.
\]
Thus for arbitrary $c\in C$, there exists $a\in A$ such that $(g\circ f)(a)=c$. Therefore $g\circ f$ is surjective.
\end{proof}

Commentary: The proof works backward from the target $c$. Surjectivity of $g$ gets you to some $b$, then surjectivity of $f$ gets you to some $a$.

\subsection{Solution: power set inclusion}

Claim: If $A\subseteq B$, then $\pow(A)\subseteq\pow(B)$.

\begin{proof}
Suppose $A\subseteq B$. Let $S\in\pow(A)$ be arbitrary. Then $S\subseteq A$. Since $A\subseteq B$, every element of $S$ is an element of $B$, so $S\subseteq B$. Therefore $S\in\pow(B)$. Hence $\pow(A)\subseteq\pow(B)$.
\end{proof}

Commentary: The elements of $\pow(A)$ are themselves sets. This is a type issue. To prove $S\in\pow(B)$, prove $S\subseteq B$.

\subsection{Solution: no smallest positive real}

\begin{proof}
Suppose, for contradiction, that $r$ is the smallest positive real number. Since $r>0$, we have $r/2>0$. Also $r/2<r$. This contradicts the assumption that $r$ is the smallest positive real number. Therefore no smallest positive real number exists.
\end{proof}

Commentary: The contradiction is exact: $r/2$ is a positive real smaller than the alleged smallest positive real.

\newpage
\section{Final Exam: From Beginner to Ready}
\chapterline

This is not a trick exam. It is a mirror. Use it to find what needs another pass.

\subsection{Part 1: Logic}

\begin{enumerate}
\item Define statement and predicate.
\item Give an example of a predicate whose truth depends on the domain.
\item Negate: $\forall x\in A,\exists y\in B,\ R(x,y)$.
\item Negate: if $n$ is prime, then $n$ is odd.
\item Explain why $\forall x\exists y$ is different from $\exists y\forall x$.
\end{enumerate}

\subsection{Part 2: Proofs}

\begin{enumerate}
\item Prove that the sum of two odd integers is even.
\item Prove that if $a\mid b$ and $a\mid c$, then $a\mid(b-c)$.
\item Prove that if $n^2$ is odd, then $n$ is odd.
\item Prove that there is no smallest positive real number.
\item Prove that $\sqrt2$ is irrational.
\end{enumerate}

\subsection{Part 3: Sets and functions}

\begin{enumerate}
\item Prove $A\setminus(B\cup C)=(A\setminus B)\cap(A\setminus C)$.
\item Prove that if $A\subseteq B$, then $\pow(A)\subseteq\pow(B)$.
\item Prove that composition of surjective functions is surjective.
\item Give a counterexample showing $f(S\cap T)=f(S)\cap f(T)$ can fail.
\item Prove $f^{-1}(S\cup T)=f^{-1}(S)\cup f^{-1}(T)$.
\end{enumerate}

\subsection{Part 4: Relations and induction}

\begin{enumerate}
\item Prove congruence modulo $n$ is symmetric.
\item Prove subset inclusion is a partial order.
\item Prove by induction that $1+3+\cdots+(2n-1)=n^2$.
\item Prove by induction that $2^n\geq n+1$ for $n\geq1$.
\item Explain why strong induction is useful for prime factorization.
\end{enumerate}

\subsection{Scoring}

If you can complete most of this with clear definitions and justified lines, start Velleman. If not, do not call it failure. You have found the exact floorboard that needs reinforcement.

\end{document}
